\subsection{对应的性质 IV：Whittaker 模型}

\ssszh 表示论的许多情形中，例如约化群的代数表示或 $(\fg,B)$-模，不可约表示都带有\nobreak 典范的生成直线，即最高权向量。

与之相对，对 $p$-adic 群 $G(\Fqt)$，大多数不可约模并没有典范的生成直线。粗略地说，原因在于 $\on{Dist}(K\backslash G/K)$ 对一切充分小的 $K$ 都非交换。

\begin{re} 对非分歧表示，存在典范的生成直线，即 $G(\Fqt)$-不变量。这条基本奇迹也就是球 Hecke 代数的交换性，在本领域至关重要。例如，它给出自守表示的张量积分解，并把非分歧自守表示的研究化归为 Hecke 特征函数与特征层的研究。\end{re}

不过，如果我们所要求的典范直线不在 $\pi$ 内部，而是 $\pi$ 的一个商，情形就好一些。确切地说，存在一个从 $G(\Fqt)\mod$ 到向量空间的典范函子，它把每个不可约表示送到一条直线或零。这就是 Whittaker 余不变量函子，可以这样理解。

为了逼近定义，不妨先考虑最高权理论一个看似直接的类比，即 $N(\Fqt)$-余不变量的 Jacquet 函子
\[
\on{Jac}: G(\Fqt)\mod \rightarrow \on{Vect}, \quad \pi \mapsto
\pi_{N(\Fqt)}.
\]
它虽然很有用，却把许多有趣的不可约表示送到零\footnote{这一现象大致就是尖点表示的存在性，即这样的单 $G(\Fqt)$-模：它不出现在抛物诱导模之内。我们要强调，这是 $p$-adic 情形以及较简单的有限李型群情形的基本事实，而约化群的代数表示论中没有这种现象。}，把另一些送到维数大于 $1$ 的向量空间。

为作修正，我们用一个通用特征标来扭转这个构造。确切地说，以 $i \in I$ 标记单根，考虑复合
\begin{equation} \label{e:whitchar}
\psi:  N(\Fqt) \rightarrow (N/[N,N])(\Fqt) \simeq \underset{i} \Pi
\hspace{.2mm}
\Fqt \xrightarrow{\on{sum}} \Fqt
\xrightarrow{\on{res}} \Fq,
\end{equation}
其中 $\on{sum}$ 是 Laurent 级数的加法，$\on{res}$ 把 Laurent 级数送到它的留数，即 $t^{-1}$ 的系数。再把它与一个适当通用的特征标 $\mathbb{F}_q \rightarrow k^\times$ 复合，就得到 $N(\Fqt)$ 的一个特征标，仍记作 $\psi$。Whittaker 余不变量由
\[
\on{Whit}: G(\Fqt)\mod \rightarrow \Vect, \quad \pi \mapsto \pi_{N(\Fqt),
\psi}
\]
给出。它把单模送到一条直线或零，并具有一个关键性质：许多尖点表示，例如 $G = GL_n$ 时的全部尖点表示，都被送到直线。

这里的一条基本性质是：$\on{Whit}$ 有右伴随，由把 $k$ 送到 Whittaker 函数空间
\[
\Fun( G(\Fqt) / N(\Fqt), \psi),
\]
给出，也就是把 $\psi$ 上诱导到 $G(\Fqt)$。

\ssszh 上述构造有如下几何类比。首先，拓扑群的同态 \eqref{e:whitchar} 提升为群归纳概形的映射
\[
\psi: LN \rightarrow L(N/[N, N]) \simeq \underset i \Pi \hspace{.2mm}
L\mathbb{G}_a \rightarrow L\mathbb{G}_a \rightarrow \mathbb{G}_a.
\]
由此，我们想从 $\mathbb{G}_a$ 拉回「通用特征标」的类比物。回忆 $G(\Fq)$ 的平凡表示与 $G$ 在 $\Vect$ 上恒等范畴作用之间的类比，我们因此需要 $\mathbb{G}_a$ 在 $\on{Vect}$ 上的一个通用作用。代数群 $H$ 在 $\Vect$ 上的作用，等同于 $H$ 上的一个特征标层，即一个 D-模 $\chi$，配上适当结合的等同
\[
\mu^! \chi \simeq \chi \boxtimes \chi,
\]
其中 $\mu: H \times H \rightarrow H$ 是乘法映射。我们要强调：这里的用法与第 \ref{ss:charsheaves} 节中较宽泛的「特征标层」不同。那里的用法是 $G(\Fq)$ 的有限维表示的特征标的类比；这里的用法是一维表示即特征标的类比。

\begin{re} 稍不正式地说，若对 $H$ 的任意闭点 $h$ 记 $\chi_h$ 为 $\chi$ 相应的 $!$-茎，则上式给出茎之间的同构
\[
\chi_{h_1 h_2} \simeq \chi_{h_1} \otimes \chi_{h_2}, \quad \text{对 } h_1,
h_2 \in H.
\]
也就是说，特征标层给出一个映入直线的同态，即映入 $\Vect^\heartsuit$ 的可逆元的同态；正如特征标给出一个映入非零标量的同态，即映入 $k$ 的可逆元的同态。
\end{re}

有了这些，$\Fq$ 的通用特征标的类比物，就是 $\mathbb{G}_a$ 上的指数 D-模 $\on{exp}(z)$。确切地说，若以 $z$ 记 $\mathbb{G}_a \simeq \mathbb{A}^1$ 上的标准坐标，这个 D-模承载指数函数 $e^z$ 所满足的微分方程，即
\[
\on{exp}(z) \simeq D_{\mathbb{G}_a} / ( D_{\mathbb{G}_a} \cdot (\partial_z -
1)),
\]
而它的乘性本质上来自恒等式 $e^{z_1 + z_2} = e^{z_1} \cdot e^{z_2}$。据此，我们以 $\psi$ 记 $LN$ 上相应的特征标层，即 $\psi^! \on{exp}(z)$，并同样可构造 Whittaker 余不变量函子
\[
\on{Whit}: \Dmod(LG)\mod \rightarrow \Vect, \quad \sC \mapsto \sC_{LN,
\psi}
\]
以及 Whittaker D-模构成的范畴
\begin{equation} \label{e:catwhitsheaves}
\Dmod(LG / LN, \psi).
\end{equation}
在转向它的谱侧对应物之前，我们想作几条帮助定位的评注。

\begin{re} 对于 $\mathbb{G}_a(k)$ 上的指数 D-模与 $\mathbb{G}_q(\Fq)$\zy{[译者按（上下文推断，无外部依据）：原文此处作 $\mathbb{G}_q(\Fq)$，按上下文应为 $\Fq$；本译本照原文排印。]} 的通用特征标之间的类比，我们来给出三种理解方式。

第一种相当朴素。为简单起见，设 $\mathbb{F}_q$ 与其素子域 $\mathbb{F}_p$ 重合。此时，若以 $e^{2\pi i / p}$ 记 $k^\times$ 中一个非平凡的 $p$ 次单位根，则通用特征标的一个具体取法是
\[
\mathbb{F}_p \rightarrow k^\times, \quad a \mapsto e^{ a\cdot 2\pi i / p },
\]
其余取法与它相差 $\mathbb{F}_p$ 中一个非零元素的倍数。也就是说，通用特征标看起来就像指数函数。对非平凡的扩域 $\Fq$，同样的说法在作必要调整后也成立，只需先用迹映射 $\Fq \rightarrow \Fp$。

第二种也相当实际。给定 $k$-向量空间 $V$，其特征标层与 $V^*$ 的闭点典范地一一对应，其中余向量 $\xi$ 对应特征标层 $\xi^! \on{exp}(z)$。特别地，当 $V = \mathbb{G}_a$ 时，在差 $\mathbb{G}_m$ 的伸缩的意义下只有一种非平凡的特征标层；正如对 $\Fq$，在差 $\Fq^\times$ 的伸缩的意义下只有一个非平凡特征标。

第三种更一般，走的是函数-层对应。首先，正特征的一个美妙之处在于：$\Fq$ 上任意交换群概形 $H$ 都有一个有趣的自同源，即 Lang 同源，其形状为
\[
1 \rightarrow H(\Fq) \rightarrow H \rightarrow H \rightarrow 1,
\]
也就是说，它以有理点构成的有限群为核。对 $H = \mathbb{G}_a$，具体就是短正合列
\[
1 \rightarrow \mathbb{G}_a(\Fq) \rightarrow \mathbb{G}_a \xrightarrow{x
\mapsto x^q - x} \mathbb{G}_a \rightarrow 1.
\]
沿此同源前推常值层并取其直和项，得到 $\mathbb{G}_a$ 上一族特征标层，以 $\mathbb{G}_a(\Fq)$ 的不可约表示为指标。取 Frobenius 的迹之后，这些特征标层恰好给出 $\mathbb{G}_a(\Fq)$ 的不可约特征标。也就是说，$\Fq$ 的特征标可以典范地提升为 $\mathbb{G}_a$ 上的特征标层。

对这些特征标层，若只把它们看作 $\mathbb{A}^1$ 上的 étale 局部系统，则可以问它们在 $\infty$ 处的行为，即紧化到 $\mathbb{P}^1$ 之后的行为。结论是：它们在这个剩余点处野分歧。于是，它们与特征零情形中 $\on{exp}(z)$ 的对应，就纳入野分歧 Galois 表示与带非正则奇点的 de Rham 局部系统之间更广泛的类比之中。
\end{re}

\begin{re} 我们来评注 Whittaker D-模范畴 $\Dmod(LG / LN, \psi)$ 的定义。它的定义是取 $\Dmod(LG)$ 的 $(LN, \psi)$-余不变量，与半无限旗流形上 D-模的范畴类似。根据 Raskin 的一个重要定理 \cite{whit}，它同时也等同于 $(LN, \psi)$-不变量。在 $p$-adic 情形中，Whittaker 不变量与余不变量并不字面地等同；但上述等同仍可类比于 Rodier \cite{rodier} 对 Whittaker 模型的紧逼近。
\end{re}

\ssszh 现在来描述 Whittaker 模型预期的谱侧对应物。它就是 $$\QCoh(\LS),$$ 即局部系统模空间上的全部拟凝聚层。

这里的基本直觉同样可以取自 $p$-adic 情形。$G(\Fqt)$ 的每个不可约表示 $\pi$，其 Whittaker 余不变量至多构成一条直线。改用
\[
\Hom( \on{ind}_{N(\Fqt)}^{G(\Fqt)} k_{-\psi}, \pi) =: \ell,
\]
会更方便，其中 $\on{ind}$ 是紧支诱导函子。它与 $\pi$ 的对偶表示的 Whittaker 余不变量等同，特别地仍是一条直线或零。

另一方面，在局部朗兰兹对应之下，人们给 $\pi$ 配上一个 Galois 表示 $\sigma$，并把 $\pi$ 想象成在 Galois 表示的空间 $\on{Gal}_{\cG}$ 中「坐落于 $\sigma$ 之上」。

在算术局部朗兰兹对应的现代表述中（即作为范畴的等价，参见下文第 \ref{ss:variants} 节），人们期望：凡有非零 Whittaker 余不变量的 $\pi$，都与一个支撑在 $\sigma$ 上的一维摩天楼层 $\mathbb{L}(\pi)$ 互换。而且，这条典范地联系于 $\pi$ 的直线，预期就正是 $\ell$。于是，综合等同
\[
\Hom( \on{ind}_{N(\Fqt)}^{G(\Fqt)} k_\psi, \pi) \simeq \ell \simeq
\Hom( \OO_{\on{Gal}_{\cG}},  \mathbb{L} \pi ),
\]
可以看出，$\mathbb{L}$ 应当把（紧支）诱导的 Whittaker 模与结构层 $\OO_{\on{Gal}_{\cG}}$ 互换。

在几何情形中，Whittaker 模已换成相应的 D-模空间 $\Dmod(LG/LN,\psi)$。\footnote{注意：由于 Raskin 对 Whittaker 不变量与余不变量的等同 \cite{whit}，诱导与紧支诱导之间的差别消失了。} 类似地，我们把结构层 $\OO_{\on{Gal}_{\cG}}$ 换成 $\QCoh(\LS)$；后者同样余表示整体截面，只不过现在表示的是范畴层的截面。

正如我们在讲到缠绕算子时很快就会讨论的，这个对应会引出几何与表示论中许多引人注目的具体猜想。不过目前，我们只用两条注来结束对图 \ref{f:table} 中这一行的讨论。

\begin{re} 在 $p$-adic 情形中，并非每个缓增不可约表示都有 Whittaker 模型；局部几何朗兰兹的预期则不同：每个缓增的 $\Dmod(LG)$ 表示都有 Whittaker 模型。基本想法如下。设 $x$ 是叠 $X$ 的一个点，$I_x$ 是其迷向群。则在 $X$ 上以概形论方式支撑于 $x$ 的拟凝聚层构成的阿贝尔范畴，等同于 $\on{Rep}(I_x)$；特别地，除非 $I_x$ 幂幺，摩天楼层不会生成它。在算术局部朗兰兹对应之下，这与 Shahidi 的一般分歧包猜想相吻合：该猜想断言，在一个缓增 $L$-包内部，\footnote{超出 $GL_n$ 的情形，从 $G(\Fqt)$ 的不可约表示到 Galois 表示的映射不是双射，而是有限到一；其纤维称作 $L$-包。} 恰有一个不可约表示有 Whittaker 模型 \cite{shahidi90}。

与之相对，$\QCoh(X)$ 的摩天楼模 $\Vect$ 会生成 $\TwoCoh(X)$ 中支撑于 $x$ 的子范畴；在这个意义上，每个缓增不可约模都是通用的。尽管如此，由于前面许诺的对朴素表述 \eqref{e:loclangnaive} 的修正所牵涉的问题，$\Dmod(LG/LN, \psi)$ 并不会成为 $\Dmod(LG)$ 的生成元。
\end{re}
